\begin{abstract}
本文描述艺术家调色 Skill v1.0.0 的实现及其现有输出记录。系统接收一张照片和可选的风格诉求，从十位艺术家的资料卡与真实作品索引中选择参考，将照片与画作共同提交给图像编辑模型，再从模型的实际返回图估计全局颜色映射，导出 $65^3$ 三维查找表，并应用于原尺寸照片。艺术家第一人称解读在最终成片之后生成，用于解释可见的色彩关系，不参与本次图像生成。该设计将参考选择、生成式配色建议、确定性 LUT 渲染与角色解释分开归档，降低不同阶段被混用的可能。本文说明 Oklab 空间中的高斯最优传输拟合、文件身份检查、保存后重读渲染，以及面向豆包等宿主的能力分支。结果部分只汇总既有的十图十版与三图六版记录；后者的三张照片与前批重叠，不构成十六个独立样本。十六份记录均通过所保存 LUT 与导出 PNG 的采样回放核对，但这一检查不代表 LUT 与生成图逐像素一致，也不证明审美最佳或题材泛化能力。本次发布不增加独立试跑。
\end{abstract}

\begin{IEEEkeywords}
参考图调色，三维查找表，Oklab，最优传输，艺术家角色，Agent 工作流。
\end{IEEEkeywords}

\section{引言}
摄影调色通常需要同时回答两个问题：希望图像呈现什么样的色彩关系，以及如何把这种关系落实为能够重复应用的变换。艺术家的作品可以提供具体的视觉参考，但艺术史中的个人风格并不是固定的 RGB 配方。相同艺术家在不同阶段、题材与创作目标下，可能使用显著不同的明度、色度与冷暖组织方式。因此，仅向图像模型输入一个艺术家姓名，难以说明最终颜色来自哪件作品，也难以核对其与该艺术家实际实践之间的关系。

本项目采用真实作品作为视觉依据，以艺术家资料卡辅助选图，再以可保存的 LUT 承载最终变换。它服务于摄影风格化与交流演示：用户可以自动匹配艺术家，也可以指定两位艺术家，让同一张照片得到各自独立的版本。最终交付保留相机照片的像素位置和尺寸；生成模型提供配色样稿，样稿中的新增笔触、物体及形状不直接替换相机原片。

本文中的“角色蒸馏”指把作品、创作阶段、相关生平与有来源的审美主张组织成可调用的资料卡和提示指令。它不表示训练出了模仿艺术家的神经网络，也不表示进行了参数蒸馏、强化学习或人格复原。当前版本的技术贡献是工作流与交付契约：将艺术家选择、真实画作输入、生成结果归档、LUT 计算和成片解释明确连接起来。底层颜色映射沿用高斯近似下的闭式最优传输思想\cite{pitie}，本文不将其作为新提出的颜色传输算法。

技术报告使用 IEEEtran 的 journal 双栏结构排版，模板依据 IEEE 的期刊模板说明\cite{ieee}。本文是当前软件版本的技术记录，并非已经过同行评审的期刊论文。开源入口与安装说明由项目仓库维护\cite{repo}。

\section{任务定义与系统结构}
\subsection{输入、输出与角色分工}
记用户原片为 $I$，用户的可选目标为 $u$，艺术家候选集合为 $\mathcal{A}$。一次完整任务依次选择艺术家 $a\in\mathcal{A}$ 和一幅有来源的作品 $R_a$，调用宿主图像编辑器 $G$ 得到样稿 $J_a$，再估计全局颜色函数 $f_a$。最终照片为
\begin{equation}
  O_a(p)=\widehat f_a\bigl(I(p)\bigr),
  \label{eq:render}
\end{equation}
其中 $p$ 是原片像素位置，$\widehat f_a$ 是从保存的三维 LUT 读取并插值得到的函数。输出位置由 $I$ 决定，不使用 $J_a$ 中重绘后的几何结构。

系统把四种图像当作不同对象：原片 $I$、真实参考作品 $R_a$、图像模型返回 $J_a$ 和 LUT 成片 $O_a$。参考作品决定配色来源；模型返回是拟合目标；LUT 成片是最终摄影结果。不能把模型返回称为画家原作，也不能把直接生成结果作为已经应用 LUT 的照片交付。这一区分同样适用于网页对照、视频演示和素材命名。

\begin{figure}[t]
\centering
\fbox{\begin{minipage}{0.93\columnwidth}
\small
原片与用户目标\par
$\downarrow$\par
既有艺术家路由 $\rightarrow$ 选择一幅真实作品\par
$\downarrow$\par
原片小图 + 作品小图 $\rightarrow$ 同次双图编辑\par
$\downarrow$\par
实际返回文件 $\rightarrow$ Oklab 统计拟合\par
$\downarrow$\par
保存 LUT $\rightarrow$ 重读 LUT $\rightarrow$ 原尺寸照片\par
$\downarrow$\par
根据四图及拟合记录，生成艺术家后置解读
\end{minipage}}
\caption{默认执行顺序。人物语言位于渲染之后，没有指向本次图像生成的反向控制边。}
\label{fig:workflow}
\end{figure}

\subsection{作品依据先于人格解释}
图像生成阶段使用短提示，明确两张图的角色：第一张是内容照片，第二张是色彩参考，要求把第一张的颜色向第二张匹配。当前默认英文提示为：
\begin{quote}
\small Image 1 is the content photo. Image 2 is the color reference. Match the colors of Image 1 to Image 2.
\end{quote}
提示中不加入艺术家模拟对话，也不把后续解释预先写成一长串局部操作。两张实际输入与提示原文一同保存，作品作为第二输入的事实由文件记录支持，而不是仅凭回复中提到画作名称。

图像编辑由宿主提供。宿主实际工具名称、接受的文件形式和返回格式可能不同，Skill 不把 Codex 特定参数视为豆包或其他平台的固定接口。当前流程默认准备最长边为 720 像素的输入，分别保持原片与作品的宽高比；工具的实际返回尺寸按真实文件记录，不把提示中的尺寸愿望当作已执行的输出规格。

\subsection{可重复性的范围}
本系统追求的是可追溯、可恢复的算法步骤。给定归档后的相同文件，局部统计、LUT 导出及重读渲染具有明确计算路径。艺术家选择仍由宿主模型按照规则判断，图像编辑器也可能含有不可控采样。因此，固定文字提示不保证下一次生成相同像素，固定路由规则也不等于训练得到具有校准置信度的分类器。

任务记录必须保存实际选中的角色、作品及来源、输入哈希、上传小图、原始提示、返回文件、LUT 和成片。重复调用图像编辑器需要新建版本目录；不能覆盖原返回后继续沿用旧的拟合说明。恢复执行时应从已存在且身份一致的文件继续，而不是用一个看起来相似的图片替代原参考或实际返回。

\section{艺术家路由与作品库}
\subsection{从图像关系选择艺术家}
当前候选范围包含莫奈、梵高、马蒂斯、高更、霍普、罗斯科、莫兰迪、塞尚、夏加尔和安塞尔·亚当斯。路由首先遵守用户明确指定的艺术家或参考作品；自动模式再考虑照片中已存在的光、颜色及明暗关系，并结合用户希望的效果。题材标签仅提供背景：同一座建筑可以因日照与阴影适合霍普，也可以因大面积深浅色场适合罗斯科。

路由记录应指出可定位的图像证据，例如“黄叶与冷蓝水面相邻”“受光墙面与棚下暗部可分”“白色花瓣与背景绿色接近”。随后说明这处关系为何适合选中艺术家的某项实践。仅写“这是风景，所以选择莫奈”不足以支持参考图选择。陌生题材仍在既有候选中寻找可迁移的颜色关系，不自动扩展为持续搜索新艺术家的任务。

\begin{table}[t]
\caption{路由观察线索示例。它们是选图入口，不是固定滤镜定义。}
\label{tab:routing}
\centering\small
\begin{tabular}{p{0.19\columnwidth}p{0.72\columnwidth}}
\hline
艺术家 & 可观察的关系与比较重点\\\hline
莫奈 & 相邻冷暖、环境光下的颜色连续性\\
梵高 & 既有色组对照及并存的情绪诉求\\
马蒂斯 & 多个可分色组与相互呼应的分量\\
高更 & 大主色面、简化的颜色组织与主观表达\\
霍普 & 受光面、阴面及两者构成的空间区隔\\
罗斯科 & 大面积色场及其内部的深浅关系\\
莫兰迪 & 相近色与灰阶之间的细小间隔\\
塞尚 & 邻接冷暖对体积、转面与空间的作用\\
夏加尔 & 日常场景中可承载记忆联想的配色\\
亚当斯 & 以中性黑白分辨暗部、中间调和高光\\\hline
\end{tabular}
\end{table}

多位角色均适合时，规则应让选择收敛，而非无限比较。可依次比较用户目标是否吻合、原片颜色关系是否清楚、参考作品能否支持这处关系，再按稳定的角色顺序处理完全相同的候选。这里的顺序用于减少重复决策，不表示艺术家之间存在普遍高低，也不能据此声称自动找到“最佳审美”。最终效果由用户判断。

\subsection{艺术家卡与实际作品的时期关系}
每位艺术家的资料卡选取主要时期作为理解中心，但作品清单可以包括早期、成熟期及对照作品。一次任务选到不同时期的作品时，说明应使用该作品的真实年份和语境。卡片名称中的“晚期”不能覆盖具体作品的事实。例如，本次田野比较中，莫奈参考的是 1888 年安提布的《园丁之家》，不能将这一例解释为晚年睡莲的直接配色\cite{monet}。

同一艺术家内部的差异也是路由信息。梵高并非始终追求同一类高反差，高更并非只使用一种浓色，罗斯科也不能被统一归纳为棕灰。艺术家卡提供判断问题，实际参考作品提供本次视觉目标。叙述两位艺术家的区别时，应落在所选作品与最终成片上，避免为了戏剧性把艺术史组织成互斥口号。

亚当斯路线另有明确的输出模式。当前实现将拟合后的颜色再映射为三个通道相同的中性值，并以黑白作品为参考。这使黑白成为可检查的实现条件，而不是仅靠人物文字描述。该约定说明当前 Skill 的亚当斯版本，不声称一张 LUT 能涵盖其摄影、底片曝光和暗房制作的完整过程。

\subsection{公开素材与按需引用}
作品索引包含十位艺术家的 80 个参考条目。条目记录作品名、年份、出处、图像链接、研究复制图的尺寸与哈希等信息。索引的存在不表示每张复制图均可重新分发。公开包只随附已确认可分发的三张 CC0 莫奈图像，其开放依据来自克利夫兰艺术博物馆的 Open Access 政策\cite{cma}；其余作品保留元数据和来源，按宿主可用能力取得或由用户导入。

运行时取得的作品属于本地参考缓存，不应随日志或安装包自动再次发布。来源为 PDF、包含多幅作品的页面或需要裁切的图像时，应确认实际输入是目标作品的正确区域。无法取得清楚、可信的真实作品时，任务停在选图阶段，不能用艺术家肖像、其他作者的画或旧的生成样稿补齐双图输入。

\section{从生成样稿到全局 LUT}
\subsection{颜色规范化与统计估计}
拟合前读取照片及实际返回图，处理 EXIF 方向，并根据嵌入的 ICC 配置转换到 sRGB；没有配置时按 sRGB 解释并记录这一假设。当前普通摄影路径基于 SDR sRGB。输出 LUT 的输入和输出均为全范围的编码 sRGB，不能作为 Log 或 HDR 的技术转换直接使用。

颜色样本经 sRGB 解码与 Oklab 变换得到三维向量 $x=\phi(c)$。Oklab 用于组织亮度与两个对立色维度\cite{oklab}。它不是对用户情绪或艺术审美的测量器；在本实现中，它仅是统计颜色变换的计算空间。原片小图和返回图分别统计，不要求同一像素位置对应，因此即使返回图有结构改变，也可以计算整体颜色分布的近似。

对原片和目标样稿分别估计均值与协方差，记为 $(\mu_s,\Sigma_s)$ 和 $(\mu_t,\Sigma_t)$。实现采用确定性的步长抽样，目标样本量约为 $8\times10^4$；这是采样目标而非所有尺寸下恰好相同的样本数。协方差对称化并加上 $10^{-6}I_3$ 的正则项，以减小近退化颜色分布造成的数值问题。

\subsection{高斯最优传输映射}
在高斯近似下，源颜色到目标颜色的闭式映射可写为\cite{pitie}
\begin{align}
 A &= \Sigma_s^{-1/2}
       \bigl(\Sigma_s^{1/2}\Sigma_t\Sigma_s^{1/2}\bigr)^{1/2}
       \Sigma_s^{-1/2},\label{eq:transport}\\
 b &= \mu_t-A\mu_s,\\
 T(x)&=Ax+b.
\end{align}
实现通过对称矩阵特征分解计算矩阵幂，对逆平方根中的很小特征值设置数值下限，并检查 $A\Sigma_sA^\top$ 与 $\Sigma_t$ 的相对差异。该检查针对协方差计算的数值一致性，不是最终视觉相似度评分。

脚本支持多种聚合方式，本 Skill 默认的一张原片与一张实际样稿使用 \texttt{bures} 入口。当每侧只有一张图时，该入口的单元素聚合直接返回自身统计量，不存在额外的多图重心搜索。因而，本报告中的十六份输出不能被描述为比较了多种训练后端，或进行了多图联合优化。

设强度为 $\alpha$，亮度与色度分量系数为对角矩阵 $D$，映射在 Oklab 空间作用后再返回 sRGB：
\begin{align}
 z(c)&=\phi(c)+\alpha D\{T(\phi(c))-\phi(c)\},\\
 f(c)&=\operatorname{clip}_{[0,1]}\left[\phi^{-1}\bigl(z(c)\bigr)\right].
 \label{eq:mapping}
\end{align}
当前作品双图任务采用完整强度，本文所述发布整理没有为局部浓色自动降低 $\alpha$。映射的裁切属于 sRGB 输出范围处理；它可能改变部分颜色层次，但强饱和或明显色偏本身不作为任务失败判据。用户允许以强风格化方式表达配色，并保留最终审美选择权。

\subsection{离散化、保存与原尺寸渲染}
程序在 $[0,1]^3$ 中建立每维 65 个节点的规则 RGB 网格，共 $65^3=274625$ 个节点。对每个节点计算式\eqref{eq:mapping}，按 \texttt{.cube} 顺序保存，声明输入域最小值为 $(0,0,0)$、最大值为 $(1,1,1)$。读取器核对维数、节点数量、有限性及数值范围，再将结果还原为三维表。

最终照片使用从磁盘重读的表，而非直接使用内存中的连续映射。每个原片 RGB 坐标在表中通过三线性插值取得输出值，位置不变；普通 8 位路径按 128 行分块处理，并保存带 sRGB 配置的 PNG。这里保留的是原片的空间结构和输出尺寸，颜色数值会按 LUT 改变，不能将其称为无损还原或保证所有细节对比度完全不变。

保存后，脚本用固定随机种子抽取 2048 个原片位置，分别计算重读 LUT 的预期量化值和已保存 PNG 中的实际值，记录最大通道差。该抽样回放检查能发现导出错位、量化或读取路径上的不一致。它不是对所有像素重新独立证明，也不是将 $O_a$ 与 $J_a$ 比较；回放误差为零不能被改写成“与 Image 的误差为零”。

\subsection{生成重绘与全局映射的边界}
全局 LUT 对相同输入 RGB 给出相同输出 RGB，不能依据像素属于天空、叶子还是水面而分别赋色。当照片中的天空和水面具有相近蓝色，而图像编辑器只把天空改为朱红时，单一 LUT 无法完整复现这种按语义分区的操作。类似地，Image 可以把原本颜色相近的两片叶子分别染红和染绿，LUT 只能利用它们原先已有的色值差异。

当前几何筛查比较画幅比例与亮度梯度的相关性，默认要求相对画幅比例误差不超过 $0.02$，边缘相关性至少为 $0.80$。这只是辅助筛查；色彩变化也可能影响亮度梯度，阈值未在独立标注集上校准。未通过筛查时，程序仍可使用整体颜色统计拟合，但把范围记录为“几何筛查未通过后的整体统计迁色”。通过筛查同样不代表图像完全没有重绘。

全局统计路径是当前算法本身的边界，不是新加的局部修复模型。它不估计光流、不执行物体分割、不把画作笔触叠加到照片，也不把采样回放核对当作图像模型质量保证。重绘差异应在四图对照中保留，解释文本据实际 LUT 成片作答。

\section{面向不同 Agent 宿主的执行契约}
\subsection{能力分支与阶段状态}
较弱 Agent 容易把自然语言目标直接写成“已完成”，或在没有真实文件时继续描述后续结果。当前整理用明确的阶段输入和产物减少此类跳步。阶段状态由文件是否存在及其记录决定，而不是由回答是否完整决定。表\ref{tab:capability} 给出主要分支；它描述工作流要求，不代表任一宿主已在本轮实机通过。

\begin{table}[t]
\caption{宿主能力不足时的真实交付范围。}
\label{tab:capability}
\centering\small
\begin{tabular}{p{0.39\columnwidth}p{0.52\columnwidth}}
\hline
可用条件 & 可继续的阶段及交付\\\hline
能读取原片与作品 & 完成有来源的选图记录\\
支持同次双图编辑 & 提交两张实际图并保存返回\\
返回为可访问文件 & 归档返回、检查身份与尺寸\\
有 Python 及依赖 & 运行拟合、保存 LUT、渲染原片\\
已有验证后的成片 & 生成第一人称后置解读\\\hline
\end{tabular}
\end{table}

若宿主只有单图编辑能力，不能声称执行了真实画作双图追色；若无法取得编辑器实际返回文件，不能拿屏幕说明或作品本身代替拟合目标；若不能运行 Python，不能把一张生成图片命名为 LUT 成片。遇到缺失能力时，保留已完成的选择与文件，说明停在哪一步和缺少什么。后续宿主接手时可继续使用这些记录。

文件准备入口写入 \texttt{job.json}，记录原片、归档作品、上传图、提示和哈希。完成入口核对这些文件的身份，拒绝把原片、作品或上传小图误当成实际 Image 返回。若同一目录已归档不同的返回，要求建立新目录而不是覆盖。拟合完成产生 \texttt{result.json}，关联样稿、LUT、原尺寸成片及人物解读所需输入。操作说明只要求宿主使用自己真实可用的工具，不为未知平台编造函数名。

\subsection{豆包适配与比较任务}
豆包版本把安装、附原片、读取作品、双图编辑、文件保存和 Python 执行分别说明，并提示先确认当次账号和环境可用能力。安装一个包含脚本的 ZIP 并不自动证明平台具备脚本执行权限，也不证明图像返回能传入代码工具。因此，文档中的步骤和命令属于明确的执行路径；本次按用户要求没有新增独立试跑，不能据此声称已完成豆包新版的端到端兼容验证。

比较任务保持内容图身份不变，为每位指定艺术家独立选择真实参考、调用编辑器和计算 LUT。不同角色不共享同一份生成返回，也不为强化口播对比而在某一分支追加额外限制。两份 LUT 均完成后，再组织双方解释。这样可以把同图差异定位到参考选择、实际编辑返回和对应全局映射，而不把两位人物的修辞误当作生成条件。

\subsection{后置解释的内容与边界}
解释采用艺术家姓名作标签、第一人称作回答，先指出可定位的颜色关系，再连接有来源的创作实践，最后说明本次成片怎样体现这种关注。整体说明注明这是基于作品和资料的角色模拟，避免把后写台词伪装成艺术家历史原话。默认不使用冗长隐喻，不声称进行了记录中不存在的局部手工操作。

后置解释有助于用户理解配色，但它不是算法因果证明。模型生成了哪种变化，应先由四图与文件确定；艺术家角色随后可以说明这项可见关系与哪种实践相近。若最终 LUT 没有复制生成样稿中的深蓝天空，角色不能用一段诗意台词把深蓝描述成已存在。相反，用户接受高饱和、强色偏或局部浓色时，解释也不应自动把这些风格选择改写成需要纠正的问题。

\section{既有输出记录与可核对结果}
\subsection{记录范围}
本文结果来自两批已经完成的任务：第一批为十张原片对应十位艺术家的十版输出，第二批为三张原片分别交给两位艺术家的六版输出。第二批三张照片与前批重叠。因此，十六版是十六次输出记录，不能称为十六个独立测试样本；这些记录也不是训练、验证、测试严格隔离的标准评测。

第一批覆盖人物、静物、花卉、建筑、海景与风光，第二批用于同图对比。题材覆盖的存在可以展示当前工作流的用法，却不能证明路由在未知题材上普遍准确，更不能证明某位艺术家对相应摄影题材“最好”。本报告没有以这两批记录训练新的路由器，也没有新增交叉验证、用户盲评或显著性检验。

\begin{table*}[t]
\caption{两份既有交付审计的汇总。回放核对针对保存 LUT 与导出照片，不针对 Image 与 LUT 的相似度。}
\label{tab:audit}
\centering\small
\begin{tabular}{lccccc}
\hline
记录批次 & 批内原片数 & 输出版数 & $65^3$ LUT & 采样回放最大差为 0 & Image 几何筛查通过\\\hline
十位艺术家样张 & 10 & 10 & 10/10 & 10/10 & 1/10\\
三组同图比较 & 3 & 6 & 6/6 & 6/6 & 0/6\\
输出记录合计 & 不作为独立样本数相加 & 16 & 16/16 & 16/16 & 1/16\\\hline
\end{tabular}
\end{table*}

两份审计均记录最终图尺寸与原片转正后尺寸一致，并核对源文件、交付副本和保存的 LUT。第一批最小的示例尺寸包括 $3872\times2592$，最大为 $11656\times8742$；第二批三组最终尺寸分别为 $8272\times6200$、$4896\times3264$ 和 $11656\times8742$。这些数值说明实际导出尺寸，不表示图像模型在同样尺寸下生成，也不是速度或显存基准。

\subsection{回放一致性与几何筛查}
十六份记录的保存 LUT 重读渲染均通过已有的 8 位采样回放核对，最大差记录为 0。原尺寸文件确实由对应 LUT 与原片生成，是此项核对支持的结论。该结果不涉及另一生成种子、另一宿主或另一平台的重复执行，也不能作为整条生成式流程完全确定的证据。

几何筛查只有第一批亚当斯例通过，其余十五份记录均标注为几何筛查未通过后的整体颜色统计拟合。这一比例与样稿中可见的云、叶片、纹理或建筑细节重绘相符，但不能直接解释为十五份最终照片几何错误。最终照片依据式\eqref{eq:render} 使用原片像素位置；筛查评估的是 Image 返回与原片之间的结构接近程度。

亚当斯例的 LUT 输出通道差为 0，符合该路线的中性黑白设定。其余彩色路线的强度与颜色分离不同，审计不对审美给出自动通过分。用户是否喜欢某一版、是否愿意用于视频或日常照片，仍需人工选择；强烈色彩也不因偏离自然色就自动被否决。

\subsection{三组同图比较的实际差异}
莫奈与梵高使用同一张层叠田野照片。莫奈参考 1888 年《园丁之家》\cite{monet}，梵高参考 1890 年《雷雨云下的麦田》\cite{vangogh}。现有成片记录中，莫奈版偏浅青绿、黄色及淡粉紫，梵高版黄绿与青蓝的对置更明显。梵高的 Image 样稿包含较深的蓝色天空及绘画式云层，但最终 LUT 天空更浅、偏青绿，不能声称已复制参考画的雷雨天空。

马蒂斯与高更使用同一张秋叶与远水照片。马蒂斯参考 1953 年《蜗牛》\cite{matisse}，高更参考 1888 年《布道后的幻象》\cite{gauguin}。现有 LUT 成片分别呈现蓝紫水天与多色叶片，以及橙红云层、红叶和暗色水面。两版均可保持鲜艳，区别主要在多组颜色的相互呼应与红橙主调的集中，不能简单解释为一位追求鲜艳、另一位更追求鲜艳。

这组还展示全局映射的具体边界。马蒂斯 Image 返回按叶片位置分别赋色，LUT 只能依据原先的颜色差异产生多色分布；高更 Image 返回把天空整体改为朱红、同时保留蓝色水面，而 LUT 得到的天空、云层和水面关系不同。作品、Image 和最终 LUT 必须分别展示，这些实际差异不能通过人物台词抹平。

霍普与罗斯科使用同一座红瓦车站。霍普参考 1930 年《清晨的星期天》\cite{hopper}，罗斯科参考 1969 年《无题（棕与灰）》\cite{rothko}。成片中，霍普版保留更鲜明的蓝青、绿灰与暖墙、红棚之间的区隔；罗斯科版将颜色归并到较暗的棕灰范围。这可以说明本次一版偏建筑光与空间，一版偏深浅色场的组织，却不能把“冷清”和“沉静”定义为两位艺术家的固定心理标签。

\subsection{本报告没有报告的指标}
当前审计没有给出这十六份记录的 PSNR、SSIM 或四项损失汇总，本报告不补造数值。代码中的辅助损失包含 RGB 绝对差、CIEDE2000 色差、亮度 DSSIM 与 Oklab sliced Wasserstein 距离，但模块存在不等于本轮已运行该评估。本文也不把此前其他版本的分数或提示词搜索曲线移植到当前版本。

当生成模型改变结构时，逐像素指标会同时反映错位、纹理变化和颜色变化。SSIM 或 RGB 差可以辅助描述某些落差，却不适合单独承担“调色是否符合艺术家理念”的判据。分布距离能减弱位置敏感性，但不能判断一个颜色是否出现在用户期望的物体上。后续若设计定量评测，应预先区分渲染正确性、生成与 LUT 的接近程度、风格区分度及人工审美，分别报告。

\section{局限与后续工作}
当前路由是有资料约束的规则化语言判断，没有从独立标注数据中学习最优选择，也没有经校准的成功概率。80 个作品条目是可追溯的参考集合，不是能够统计代表艺术家全部创作的大样本数据库。作品的主要颜色受数字复制、裁切、光照与色彩管理影响；参考图版的平均颜色不能等同于原作材料或展厅中的观看经验。

高斯拟合只使用一阶和二阶统计，能够产生平滑、可导出的整体颜色映射，但会忽略复杂的多峰分布、物体语义与空间关系。图像编辑器的局部重配色不能普遍由一张 LUT 表达。未来可以在保持原片结构和明确交付边界的前提下，比较更丰富的颜色分布模型或加入显式的局部方案；这属于后续设计，不是当前版本已经具有的能力。

宿主工具的可用性和返回形态仍影响整条流程。当前豆包适配提供明确说明、能力分支和恢复条件，本次按用户要求直接整理发布，没有新增端到端实机验证。宿主不能提交双图、不能保存返回或不能执行 Python 时，必须如实交付当前阶段，而不能依靠较完整的语言说明弥补缺失的图像或 LUT。

后置角色解释用于交流，不能证明艺术家若面对这张照片必然会作同样选择，也不能替代用户评估。未来的系统评测需要明确样本划分、同图多角色设计、真实独立的用户判断，以及生成调用的实际成本与随机性记录。对于强风格化任务，预先约定哪些变化是用户接受的审美表达，比用统一的“自然”阈值自动压低色彩更符合本系统目标。

\section{结论}
艺术家调色 Skill 将真实作品参考、双图编辑、全局 LUT 和成片后角色解释组织成一条可检查的流程。当前版本强调三项区分：画作提供本次视觉来源，Image 返回提供实际拟合目标，最终照片来自原片经保存 LUT 的全局映射；人物语言在这些步骤之后生成。针对较弱 Agent 的优化集中在明确输入、文件身份、成功条件和能力不足时的真实交付范围。

既有十六份输出记录支持了保存 LUT 与原尺寸导出照片之间的采样一致性，也清楚暴露了生成重绘与全局颜色函数之间的落差。它们不支持“十六个独立测试样本”“最佳艺术家选择”或跨宿主普遍通过等结论。当前公开版本保留可用的风格强度与人工审美选择权，在这个边界内提供可归档、可交接的摄影调色工具。
